\subsubsection{常见博弈游戏}
\begin{multicols}{2}
%\paragraph{减法游戏}$n$根火柴每次可以拿$a_i$根, 最后一个拿的人获胜.
%\paragraph{Ferguson游戏}两个盒子, 一个有$m$个糖另一个有$n$个糖, 每次走步将一个盒子清空而把另一个盒子的一些糖拿到被清空的盒子中, 使得两个盒子各至少有一颗糖. 到终态(1,1)的人获胜.
%\paragraph{动态减法游戏}初始有$n$根火柴, 第一个人可以任意取, 至少$1$根至多$n-1$根, 以后每个人至少拿一根, 但拿的火柴数不能超过上一次拿的火柴数.
\paragraph{Nim-K游戏}$n$堆石子轮流拿, 每次最多可以拿k堆石子, 谁走最后一步输。结论: 把每一堆石子的sg值(即石子数量)二进制分解,先手必败当且仅当每一位二进制位上1的个数是(k+1)的倍数.
\paragraph{Anti-Nim游戏}$n$堆石子轮流拿, 谁走最后一步输。结论: 先手胜当且仅当1. 所有堆石子数都为1且游戏的SG值为0 (即有偶数个孤单堆-每堆只有1个石子数) 2. 存在某堆石子数大于1且游戏的SG值不为0.
\paragraph{斐波那契博弈}有一堆物品, 两人轮流取物品, 先手最少取一个, 至多无上限, 但不能把物品取完, 之后每次取的物品数不能超过上一次取的物品数的二倍且至少为一件, 取走最后一件物品的人获胜. 结论: 先手胜当且仅当物品数$n$不是斐波那契数.
\paragraph{威佐夫博弈}有两堆石子, 博弈双方每次可以取一堆石子中的任意个, 不能不取, 或者取两堆石子中的相同个. 先取完者赢. 结论: 求出两堆石子$A$和$B$的差值$C$, 如果$\left\lfloor C*\frac{\sqrt{5}+1}{2}\right\rfloor=min(A,B)$那么后手赢, 否则先手赢.
\paragraph{约瑟夫环}$n$ 个人 $0,\dots,n-1$, 令 $f_{i,m}$ 为 $i$ 个人报 $m$ 的胜利者, 则 $f_{1,m}=0,f_{i,m}=(f_{i-1,m}+m)\bmod i$.
\paragraph{阶梯Nim} 在一个阶梯上, 每次选一个台阶上任意个式子移到下一个台阶上, 不可移动者输. 结论: SG值等于奇数层台阶上石子数的异或和. 对于树形结构也适用, 奇数层节点上所有石子数异或起来即可.
\paragraph{图上博弈} 给定无向图, 先手从某点开始走, 只能走相邻且未走过的点, 无法移动者输. 对该图求最大匹配, 若某个点不一定在最大匹配中则先手必败, 否则先手必胜.
% 博弈论内容由 sections/RandomizationAndGameTheory/index.tex 统一引用。
\end{multicols}

\subsubsection{概率相关}
$D(X)=E(X-E(X))^2=E(X^2)-(E(X))^2,D(X+Y)=D(X)+D(Y),$\\
$D(aX)=a^2D(X), E[x]=\sum_{i=1}^{\infty}P(X\geq i)$, $m$ 个数的方差：$\begin{aligned}s^2=&\frac{\sum_{i=1}^m x_i^2}m-\overline x^2\end{aligned}$
